\documentclass[12pt,a4paper]{article}
\usepackage[utf8]{inputenc}
\usepackage[T1]{fontenc}
\usepackage{geometry}
\geometry{margin=1in}
\usepackage{hyperref}
\usepackage{listings}
\usepackage{xcolor}
\usepackage{graphicx}
\usepackage{booktabs}
\usepackage{titlesec}

% Define SmartShop Luxury Theme Colors
\definecolor{brandblue}{RGB}{30, 64, 175}
\definecolor{brandgreen}{RGB}{22, 101, 52}
\definecolor{codegray}{rgb}{0.5,0.5,0.5}
\definecolor{backcolour}{rgb}{0.95,0.95,0.92}

\hypersetup{
    colorlinks=true,
    linkcolor=brandblue,
    filecolor=magenta,      
    urlcolor=brandblue,
    pdftitle={SmartShop Technical Documentation},
}

\title{
    \textbf{SmartShop: High-Fidelity E-Commerce Ecosystem} \\
    \large Technical Architecture, Design Philosophy, and System Hardening Report
}
\author{LUWI Mafuleti --- IRISI}
\date{May 23, 2026}

\begin{document}

\maketitle
\tableofcontents
\newpage

\section{Introduction}

The SmartShop project represents a sophisticated transition from a functional Laravel skeleton into a production-grade e-commerce architecture. Designed with a focus on \textit{Luxury Retail UX} and \textit{Deterministic System Integrity}, the platform serves as a benchmark for modern web development in the context of Computer Networks and Information Systems.

This document serves as the master technical blueprint, detailing the evolution of the platform from its initial migrations to its current state: a hardened, payment-integrated, and API-ready ecosystem.

\section{The Vision: Luxury E-Commerce Philosophy}

Unlike standard template-driven marketplaces, SmartShop adheres to a "High-Fidelity" design philosophy:
\begin{itemize}
    \item \textbf{Visual Minimalism:} Leveraging a custom Vanilla CSS variable-based theme engine to achieve a lightweight, responsive, and luxury aesthetic.
    \item \textbf{Operational Resilience:} Prioritizing data integrity through transactional logic and row-level locking.
    \item \textbf{Asynchronous Excellence:} Utilizing queued mail systems and background workers to ensure zero-latency user experiences.
\end{itemize}

\section{Technological Stack}

The platform is built on a "Hardened Modern" stack, ensuring scalability and security at every layer.

\subsection{Backend: The Laravel 13 Framework}
LUWI utilizes \textbf{Laravel 13}, taking advantage of the latest PHP 8.3 features. Key backend components include:
\begin{itemize}
    \item \textbf{Eloquent ORM:} For semantic and efficient database interactions.
    \item \textbf{Laravel Sanctum:} Providing a lightweight authentication system for both SPA/Mobile and traditional web sessions.
    \item \textbf{Database Transactions:} Ensuring atomic operations during complex checkout and cancellation flows.
\end{itemize}

\subsection{Frontend: The Blade-Vanilla Synthesis}
In a departure from heavy JavaScript frameworks, LUWI employs:
\begin{itemize}
    \item \textbf{Blade Templating:} For high-performance server-side rendering.
    \item \textbf{Vanilla CSS Variables:} Enabling a robust Dark/Light theme engine without the overhead of utility-first frameworks.
    \item \textbf{Componentization:} Reusable Blade components (e.g., \texttt{product-card}) ensure a single source of truth for UI elements.
\end{itemize}

\subsection{Integrations: The Financial \& Support Tier}
\begin{itemize}
    \item \textbf{PayPal (\texttt{srmklive/paypal}):} Server-side payment capture and verification.
    \item \textbf{SMTP / Mail Queue:} Gmail-integrated notification system with database-backed queueing.
\end{itemize}

\section{System Architecture: A Layered MVC Approach}

LUWI follows a strict layered MVC (Model-View-Controller) architecture, further hardened with specialized validation and authorization layers.

\subsection{1. Presentation Layer (Blade Views)}
The UI is composed of structured layouts and partials, ensuring consistency across Public, Member, and Administrative interfaces.

\subsection{2. Application Layer (Controllers)}
Controllers are kept "lean," delegating data retrieval to Eloquent and validation to dedicated Request classes.
\begin{itemize}
    \item \texttt{ViewController}: Orchestrates public catalog and product discovery.
    \item \texttt{OrderController}: Manages the sensitive checkout-to-cancellation lifecycle.
    \item \texttt{PaymentController}: Handles the external handshake with the PayPal API.
\end{itemize}

\subsection{3. Domain Layer (Eloquent Models)}
Models encapsulate the business logic and relationships. LUWI implements a multi-model identity system, linking Users to Addresses and tracking financial state through a dedicated Payment model.

\subsection{4. Security \& Access Control Tier}
\begin{itemize}
    \item \textbf{Middlewares:} Restrict access to Member-only and Admin-only routes.
    \item \textbf{Authorization Policies:} Enforce row-level ownership rules (e.g., users can only view their own orders).
    \item \textbf{CSRF \& Mass Assignment Protection:} Enforced globally.
\end{itemize}

\section{Database Architecture \& Schema Design}

The database is optimized for transactional consistency and historical auditing.

\subsection{1. Core Relational Schema}
The platform's heart lies in its relational integrity, ensuring that products, orders, and payments are permanently linked.
\begin{itemize}
    \item \textbf{Users \& Roles:} Differentiates between \texttt{user} and \texttt{admin} tiers using an \texttt{is\_admin} flag.
    \item \textbf{Products \& Inventory:} Tracks stock levels, descriptions, and categories.
    \item \textbf{Cart \& Order Life-cycle:} Implements a transition from session-persistent \texttt{CartItems} to immutable \texttt{OrderItems} upon checkout.
    \item \textbf{Payments:} A dedicated ledger for PayPal transaction IDs and statuses.
\end{itemize}

\subsection{2. Logical Data Model (MLD) Summary}
\begin{itemize}
    \item \texttt{users} (1:N) \texttt{orders}
    \item \texttt{products} (1:N) \texttt{order\_items}
    \item \texttt{orders} (1:1) \texttt{payments}
    \item \texttt{users} (1:N) \texttt{addresses}
\end{itemize}

\section{Git Workflow \& Contribution Visibility}

As required by the project guidelines, the development was managed using Git to ensure version control and logical progression.

\subsection{1. Commit Strategy}
The project followed a semantic commit history:
\begin{itemize}
    \item \textbf{Chore:} Initial environment setup and framework installation.
    \item \textbf{Feat:} Incremental implementation of the database schema and backend logic.
    \item \textbf{UI:} Development of the luxury-themed frontend.
    \item \textbf{Docs:} Finalization of technical documentation.
\end{itemize}

\subsection{2. Git History Evidence}
[Insert Screenshot of Git Log here]
\textit{The repository history (available at the provided Git link) shows the chronological evolution of the system from its initial state to the hardened version.}

\section{System Implementation & Screenshots}

[Insert Application Screenshots here with descriptions]

\section{System Analysis: Pros, Cons, and Hardening}

\subsection{1. The Checkout Pipeline (Atomic Transactions)}
The \texttt{OrderController@store} method is the most sensitive operation in the system. It implements:
\begin{enumerate}
    \item \textbf{Stock Locking:} Uses \texttt{lockForUpdate()} to prevent race conditions during high-concurrency flash sales.
    \item \textbf{Price Freezing:} Captures the product's price at the moment of checkout, protecting the order from future price fluctuations.
    \item \textbf{Atomic Commit:} Executes the entire process within a \texttt{DB::transaction()} to ensure that an order is never created without stock reduction.
\end{enumerate}

\subsection{2. The Payment Handshake (PayPal)}
The integration uses a server-side redirect flow:
\begin{itemize}
    \item \textbf{Initiation:} Generates a PayPal order ID and stores it as a "Pending" payment.
    \item \textbf{Capture:} Upon return, the server verifies the transaction status directly with PayPal's API before marking the order as "Paid."
\end{itemize}

\subsection{3. Notification Engine (Queued Workers)}
LUWI implements a 4-tier notification suite:
\begin{itemize}
    \item \textbf{WelcomeMember:} Triggered upon registration.
    \item \textbf{OrderConfirmed:} Triggered after checkout.
    \item \textbf{OrderCancelled:} Triggered upon user-led cancellation (restoring stock).
    \item \textbf{PaymentSuccess:} Triggered after successful PayPal capture.
\end{itemize}

\section{System Analysis: Pros, Cons, and Hardening}

\subsection{Strengths (Pros)}
\begin{itemize}
    \item \textbf{Performance:} Server-side rendering (Blade) and Eager Loading (\texttt{with()}) provide sub-second page loads.
    \item \textbf{Security:} Enforces strict Mass Assignment protection and Authorization Policies.
    \item \textbf{Maintainability:} Highly modular architecture following standard PSR standards.
    \item \textbf{UX Clarity:} High-contrast theme engine and real-time email feedback.
\end{itemize}

\subsection{Limitations (Cons)}
\begin{itemize}
    \item \textbf{Infrastructure Dependency:} Requires a configured queue worker (\texttt{queue:work}) for notifications.
    \item \textbf{Scaling Threshold:} Current file-based image storage should be migrated to S3/Cloudinary for multi-region scaling.
\end{itemize}

\section{Current Status \& Readiness}

The system is currently in the \textbf{Pre-Deployment Readiness} phase.
\begin{itemize}
    \item \textbf{Status:} Functional. PayPal Sandbox loop resolved.
    \item \textbf{Readiness:} 98\% (Pending Redis Cache implementation for the Product Catalog).
\end{itemize}

\section{Conclusion}

LUWI is more than a store; it is a robust digital ecosystem built on the principles of modern software engineering. By combining Laravel's powerful abstractions with custom-hardened logic, it provides a secure, visually stunning, and highly performant platform ready for the luxury retail market.

\end{document}
